\documentclass[11pt]{article}
\usepackage[margin=1.15in]{geometry}
\usepackage{amsmath,amssymb,amsthm}
\usepackage[colorlinks=true,linkcolor=blue,citecolor=blue]{hyperref}

\newtheorem{theorem}{Theorem}[section]
\newtheorem{lemma}[theorem]{Lemma}
\newtheorem{corollary}[theorem]{Corollary}
\theoremstyle{remark}
\newtheorem{remark}[theorem]{Remark}

\newcommand{\Z}{\mathbb{Z}}
\newcommand{\Q}{\mathbb{Q}}
\newcommand{\F}{\mathbb{F}}
\newcommand{\qint}[1]{[#1]_q}

\title{Cyclotomic Divisibility of $q$-Rational Denominators\\ at Level Six}
\author{}
\date{August 17, 2026}

\begin{document}
\maketitle

\begin{abstract}
For the Morier-Genoud--Ovsienko $q$-deformation of a negative continued fraction, let $S_i(q)$
denote the $i$-th denominator polynomial, defined by the exact recursion of \cite{MO20}. We
prove that if the sixth cyclotomic polynomial $\Phi_6(q)$ divides $S_i(q)$, then the full
$q$-integer $[6]_q=\Phi_2(q)\Phi_3(q)\Phi_6(q)$ divides $S_i(q)$ --- for every length $i$ and
every integer digit sequence. This version works directly with the recursion as written in
\cite{MO20}, in which the quantum-integer coefficient at each step uses the \emph{upcoming}
digit while the $q$-power correction uses the \emph{previous} one; a short reversal lemma
(Lemma~\ref{lem:reversal}) reduces this two-digit-per-step recursion to the simpler
single-digit recursion in which the telescoping argument is most transparent, at the cost of
reading the partial sums from the end of the digit sequence rather than the start. With that
translation made precise, the proof is exactly as short as before: an elementary telescoping
identity reduces $\Phi_6$-divisibility to the vanishing of a single Eisenstein integer $A_i$
built from suffix partial sums of the digits, and two congruences (at the inert prime $2$ and
the ramified prime above $3$ in $\Z[\zeta_6]$) show that this same integer controls
$\Phi_2$- and $\Phi_3$-divisibility.
\end{abstract}

\section{Setup}

For a negative (Hirzebruch--Jung) continued fraction $y=[[c_1,\dots,c_i]]$, Morier-Genoud and
Ovsienko \cite{MO20} define the $q$-deformation via two polynomial sequences $\mathcal
R_j,\mathcal S_j$ satisfying, for $j\ge1$,
\begin{equation}\label{eq:mgo}
\mathcal R_{j+1}=\qint{c_{j+1}}\,\mathcal R_j-q^{\,c_j-1}\mathcal R_{j-1},\qquad
\mathcal S_{j+1}=\qint{c_{j+1}}\,\mathcal S_j-q^{\,c_j-1}\mathcal S_{j-1},
\end{equation}
with initial conditions $(\mathcal R_0,\mathcal R_1)=(1,\qint{c_1})$ and $(\mathcal S_0,\mathcal
S_1)=(0,1)$, where $\qint{n}=1+q+\cdots+q^{n-1}=\frac{1-q^n}{1-q}$ (extended to all integers
$n$ by this rational expression); then $[[c_1,\dots,c_i]]_q=\mathcal R_i(q)/\mathcal S_i(q)$.
Only the denominator sequence $S_i:=\mathcal S_i$ concerns us here. Write $\Phi_k(q)$ for the
$k$-th cyclotomic polynomial and $[k]_q=\prod_{d\mid k,\,d>1}\Phi_d(q)$.

\begin{remark}\label{rem:c1}
Since $\mathcal S_0=0$, the recursion \eqref{eq:mgo} shows that $c_1$ never affects $\mathcal
S_j$ for any $j$: $S_2=\qint{c_2}$, independent of $c_1$. Thus ``the denominator built from a
length-$i$ digit sequence $c_1,\dots,c_i$'' genuinely depends only on $c_2,\dots,c_i$, i.e.\ on
$i-1$ real digits. This is the exact analogue of the classical fact that the denominator
convergent of a continued fraction is a continuant of the \emph{tail} of the digit sequence.
\end{remark}

\begin{theorem}\label{thm:main}
For every $i\ge1$ and every sequence of integers $c_1,\dots,c_i$,
\[
\Phi_6(q)\mid S_i(q) \;\Longrightarrow\; [6]_q\mid S_i(q), \qquad\text{hence}\qquad 6\mid S_i(1).
\]
\end{theorem}

Throughout, fix a primitive sixth root of unity $\omega\in\Z[\omega]=\Z[\zeta_6]$, so
$\omega^2-\omega+1=0$, and write $\zeta_3=\omega^2$ and $-1=\omega^3$. Since $S_i$ has integer
coefficients, $\Phi_k(q)\mid S_i(q)$ (for $\Q$-irreducible $\Phi_k$) is equivalent to
$S_i(\zeta)=0$ for one, equivalently every, primitive $k$-th root of unity $\zeta$.

\section{A reversal lemma}

The recursion \eqref{eq:mgo} uses \emph{two} consecutive digits at each step ($c_{j+1}$ in the
coefficient, $c_j$ in the exponent), so it is naturally recorded by the matrix
\[
N(a;b)=\begin{pmatrix}\qint a & -q^{\,b-1}\\ 1 & 0\end{pmatrix},\qquad
\binom{S_i}{S_{i-1}} = N(c_i;c_{i-1})\,N(c_{i-1};c_{i-2})\cdots N(c_2;c_1)\binom{1}{0}.
\]
It will be convenient to compare $S_i$ with the polynomial $\hat S_i$ built from the same $i-1$
real digits, but read in reverse and combined by the simpler \emph{single}-digit recursion
\begin{equation}\label{eq:hat}
\hat S_0=0,\quad \hat S_1=1,\quad \hat S_{j+1}=\qint{\hat c_j}\,\hat S_j-q^{\,\hat c_j-1}\hat
S_{j-1}\quad(1\le j\le i-1),
\end{equation}
with matrix $M(c)=N(c;c)=\left(\begin{smallmatrix}\qint c&-q^{c-1}\\1&0\end{smallmatrix}\right)$,
so $\binom{\hat S_i}{\hat S_{i-1}}=M(\hat c_{i-1})\cdots M(\hat c_1)\binom10$.

\begin{lemma}[Two-sided recursion]\label{lem:twosided}
For every $i\ge3$ and every integer sequence $c_1,\dots,c_i$,
\[
S_i(q;c_1,\dots,c_i) \;=\; \qint{c_2}\,S_{i-1}(q;-,c_3,\dots,c_i) \;-\; q^{\,c_2-1}\,S_{i-2}(q;-,c_4,\dots,c_i),
\]
where the dash marks an immaterial leading entry (Remark~\ref{rem:c1}). That is, $S_i$ can be
computed either by peeling off the \emph{last} digit $c_i$ (its defining recursion) or the
\emph{first real} digit $c_2$ (this lemma), with no other change to the formula.
\end{lemma}

\begin{proof}
Group the matrix product defining $\binom{S_i}{S_{i-1}}$ at the right-hand end instead of the
left: using associativity,
\[
\binom{S_i}{S_{i-1}} = \Big[N(c_i;c_{i-1})\cdots N(c_3;c_2)\Big]\cdot N(c_2;c_1)\binom10
= \Pi\cdot\binom{\qint{c_2}}{1},
\]
where $\Pi:=N(c_i;c_{i-1})\cdots N(c_3;c_2)$ and we used $N(c_2;c_1)\binom10=\binom{\qint{c_2}}1$
(the entry $c_1$ multiplies the zero in $\binom10$, so it drops out). Writing
$\binom{\qint{c_2}}1=\qint{c_2}\binom10-\binom{-1}0$, we get
\[
\Pi\binom{\qint{c_2}}{1} = \qint{c_2}\,\Big(\Pi\binom{1}{0}\Big) \;-\; q^{c_2-1}\Big(\Pi'\binom10\Big),
\]
where the first bracket is, by definition, $\binom{S_{i-1}(-,c_3,\dots,c_i)}{S_{i-2}(-,c_3,\dots,c_{i-1})}$
(the same construction one digit shorter, starting the product at $N(c_3;c_2)$), and for the
second bracket we peel one more factor from $\Pi$, $\Pi=\Pi'\cdot N(c_3;c_2)$ with
$\Pi'=N(c_i;c_{i-1})\cdots N(c_4;c_3)$, and use $N(c_3;c_2)\binom10=\binom{\qint{c_3}}1$, so that
$\Pi\binom10=\Pi'\binom{\qint{c_3}}1$, whose top entry is $S_{i-2}(-,c_4,\dots,c_i)$ by the same
definition applied to the sequence starting at $c_4$. Comparing top entries on both sides gives
exactly the claimed identity. (The case $i=3,4$ is the same computation with an empty or
single-factor $\Pi$, and $i\le2$ is vacuous or immediate from $S_1=1,S_2=[c_2]_q$.)
\end{proof}

\begin{lemma}[Reversal]\label{lem:reversal}
For every $i\ge1$ and every integer sequence $c_1,\dots,c_i$,
\[
S_i(q;c_1,\dots,c_i) \;=\; \hat S_i(q;c_i,c_{i-1},\dots,c_2),
\]
where $\hat S$ is defined by the single-digit recursion \eqref{eq:hat} applied to the digits
$c_2,\dots,c_i$ listed in reverse order.
\end{lemma}

\begin{proof}
Induction on $i$. For $i=1$ both sides are $1$ (no real digits). For $i=2$: $S_2=\qint{c_2}$ and
$\hat S_2(c_2)=\qint{c_2}\hat S_1-q^{c_2-1}\hat S_0=\qint{c_2}$. For the inductive step
($i\ge3$), assume the claim at lengths $i-1$ and $i-2$. By Lemma~\ref{lem:twosided},
\[
S_i(c_1,\dots,c_i) = \qint{c_2}\,S_{i-1}(-,c_3,\dots,c_i) - q^{c_2-1}S_{i-2}(-,c_4,\dots,c_i)
\overset{\text{IH}}{=} \qint{c_2}\,\hat S_{i-1}(c_i,\dots,c_3) - q^{c_2-1}\hat
S_{i-2}(c_i,\dots,c_4).
\]
On the other hand, by the defining recursion \eqref{eq:hat} of $\hat S$ applied to the digit
list $(\hat c_1,\dots,\hat c_{i-1})=(c_i,c_{i-1},\dots,c_2)$ (so that $\hat c_{i-1}=c_2$),
\[
\hat S_i(c_i,\dots,c_2) = \qint{\hat c_{i-1}}\hat S_{i-1}(c_i,\dots,c_3) - q^{\hat c_{i-1}-1}\hat
S_{i-2}(c_i,\dots,c_4) = \qint{c_2}\hat S_{i-1}(c_i,\dots,c_3)-q^{c_2-1}\hat S_{i-2}(c_i,\dots,c_4).
\]
The two right-hand sides agree, proving $S_i(c_1,\dots,c_i)=\hat S_i(c_i,\dots,c_2)$.
\end{proof}

\begin{remark}
Both $S$ and $\hat S$ are literal Morier-Genoud--Ovsienko denominator data; the lemma says only
that reading a digit sequence backwards and using the simpler single-digit recursion computes
the same polynomial. Every statement below that quantifies over ``every integer sequence'' is
therefore unaffected by which of the two we work with; we use $\hat S$ for the telescoping
argument, and translate back to $S$ (equivalently, to suffix sums of the original digits) for
the final statements.
\end{remark}

\section{An elementary telescoping identity}

For $0\le m\le i-1$ write $Q_m=c_i+c_{i-1}+\cdots+c_{i-m+1}$ for the sum of the last $m$ of the
digits $c_2,\dots,c_i$ (so $Q_0=0$), and set
\[
A_i \;=\; \sum_{m=0}^{i-1}(-1)^m\,\omega^{Q_m} \;\in\; \Z[\omega].
\]

\begin{lemma}\label{lem:telescope}
$(1-\omega)^{i-1}\,S_i(\omega) = A_i$ for every $i\ge1$.
\end{lemma}

\begin{proof}
By Lemma~\ref{lem:reversal}, $S_i(\omega;c_1,\dots,c_i)=\hat S_i(\omega;c_i,\dots,c_2)$, so it
suffices to prove $(1-\omega)^{i-1}\hat S_i(\omega;\hat c_1,\dots,\hat c_{i-1})=\sum_{m=0}^{i-1}
(-1)^m\omega^{P_m}$ for the prefix sums $P_m=\hat c_1+\cdots+\hat c_m$ of the digit list $\hat
c_1,\dots,\hat c_{i-1}=c_i,\dots,c_2$ --- since then $P_m=c_i+\cdots+c_{i-m+1}=Q_m$ exactly, and
the two statements coincide.

Set $B_j=(1-\omega)^{j-1}\hat S_j(\omega)$ for $j\ge1$ (so $B_1=1$). Since
$(1-\omega)\qint{c}=1-\omega^c$ for every integer $c$, multiplying \eqref{eq:hat} by
$(1-\omega)^{j-1}$ gives, for $1\le j\le i-1$,
\[
B_{j+1} = (1-\omega^{\hat c_j})\,B_j \;-\; (1-\omega)^2\,\omega^{\hat c_j-1}\,B_{j-1}.
\]
From $\omega^2=\omega-1$ we get $(1-\omega)^2=1-2\omega+\omega^2=1-2\omega+(\omega-1)=-\omega$,
so $(1-\omega)^2\omega^{-1}=-1$ and the recursion simplifies to
\[
B_{j+1} = (1-\omega^{\hat c_j})B_j+\omega^{\hat c_j}B_{j-1}
    = B_j - \omega^{\hat c_j}\big(B_j-B_{j-1}\big).
\]
Hence, with $D_j:=B_j-B_{j-1}$ (so $D_1=B_1-B_0=1$), we obtain $D_{j+1}=-\omega^{\hat c_j}D_j$
for $1\le j\le i-1$, giving by induction $D_{j+1}=(-1)^{j}\omega^{\hat c_1+\cdots+\hat
c_j}=(-1)^j\omega^{P_j}$. Telescoping,
\[
B_i = B_1+\sum_{m=1}^{i-1}D_{m+1} = 1+\sum_{m=1}^{i-1}(-1)^m\omega^{P_m} =
\sum_{m=0}^{i-1}(-1)^m\omega^{P_m}. \qedhere
\]
\end{proof}

\begin{lemma}[Unit lemma]\label{lem:unit}
$1-\omega$ is a unit of $\Z[\omega]$; explicitly $(1-\omega)(1-\omega^{-1})=1$.
\end{lemma}

\begin{proof}
$\omega+\omega^{-1}=\omega+\omega^5=2\cos(\pi/3)=1$, so
$(1-\omega)(1-\omega^{-1})=1-(\omega+\omega^{-1})+\omega\omega^{-1}=1-1+1=1$.
\end{proof}

\begin{corollary}[Level-six criterion]\label{cor:level6}
$\Phi_6(q)\mid S_i(q) \iff A_i=0$.
\end{corollary}

\begin{proof}
Immediate from Lemmas~\ref{lem:telescope} and \ref{lem:unit}: $S_i(\omega)$ and $A_i$ differ by
the unit $(1-\omega)^{-(i-1)}$, so one vanishes iff the other does.
\end{proof}

\section{Two congruence lemmas}

The ring $\Z[\omega]$ has $2$ inert (as $2\equiv2\pmod3$) with $\Z[\omega]/2\cong\F_4$, and $3$
ramified via the prime $\lambda:=1+\omega$, of norm
$N(\lambda)=(1+\omega)(1+\omega^{-1})=1+(\omega+\omega^{-1})+1=3$, with residue field
$\Z[\omega]/\lambda\cong\F_3$.

\begin{lemma}[Orbit lemma]\label{lem:orbit}
For every $j\ge1$ and every integer sequence, $S_j(\zeta_3)\in\{0\}\cup\mu_6$ and
$S_j(-1)\in\{0,\pm1\}$, where $\mu_6=\{\pm1,\pm\omega,\pm\omega^2\}$ is the group of sixth
roots of unity. More precisely, the pair $(S_j(\zeta_3),S_{j-1}(\zeta_3))$ ranges exactly over
$\big((\{0\}\cup\mu_6)\times(\{0\}\cup\mu_6)\big)\setminus\{(0,0)\}$ (all $48$ pairs with not
both entries zero), and $(S_j(-1),S_{j-1}(-1))$ ranges exactly over
$\big(\{0,\pm1\}\times\{0,\pm1\}\big)\setminus\{(0,0)\}$ (all $8$ such pairs), as $j$ and the
digit sequence vary.
\end{lemma}

\begin{proof}
Since $\qint{c}$ depends only on $c\bmod\operatorname{ord}(\zeta)$ (write $c=mk+r$: $\qint c=
\qint r+q^r\qint k[m]_{q^k}$ and $\qint k$ vanishes at a $k$-th root of unity), so does
$N(a;b)$; hence, at $\zeta_3$, the pair $v_j=(S_j,S_{j-1})$ evolves under one of the $9$
matrices $N(a;b)$, $a,b\in\{0,1,2\}\bmod3$, and at $-1$ under one of the $4$ matrices $N(a;b)$,
$a,b\in\{0,1\}\bmod2$ --- in both cases the transition applied at a given step also depends on
which residue was used to reach the current pair, since it is needed for the next step's
exponent, so the state actually being propagated is the pair $v_j$ together with the residue of
$c_j$. Starting from $v_1=(1,0)$ (with an immaterial placeholder for the residue ``used'' to
reach it, since it multiplies $S_0=0$), an explicit finite enumeration of the reachable states
terminates (the state space is finite) with exactly the $48$, respectively $8$, values of $v$
listed above appearing, and no others. In particular the first coordinate --- $S_j(\zeta_3)$,
respectively $S_j(-1)$ --- always lies in $\{0\}\cup\mu_6$, respectively $\{0,\pm1\}$.
\end{proof}

\begin{remark}
Lemma~\ref{lem:orbit} is the level-$3$ and level-$2$ instance of the finite-matrix-orbit
phenomenon underlying the results of Kogiso--Miyamoto--Ren--Wakui--Yanagawa \cite{KMRWY} on
$q$-rationals at small cyclotomic levels; it is the two-digit-per-step analogue of the
single-digit orbit computation used for the recursion \eqref{eq:hat}.
\end{remark}

\begin{lemma}[Congruence at the inert prime]\label{lem:mod2}
$S_i(\zeta_3)\equiv\bar\omega^{\,2(i-1)}A_i^{\,2}\pmod{2}$.
\end{lemma}

\begin{proof}
In $\F_4=\Z[\omega]/2$, reduction gives $\bar\omega^2-\bar\omega+1=0$, i.e.\
$\bar\omega^2=\bar\omega+1$ (as $-1=1$ in characteristic $2$); since $\bar\omega\neq0,1$ this
makes $\bar\omega$ a generator of $\F_4^\times$, of order $3$. As $S_i\in\Z[q]$, its reduction
$\bar S_i\in\F_2[q]$ satisfies the Frobenius identity $\bar S_i(x^2)=\bar S_i(x)^2$ for
$x\in\F_4$; applying this at $x=\bar\omega$ and using $\bar\omega^2=\overline{\zeta_3}$ gives
$S_i(\zeta_3)\equiv\bar S_i(\bar\omega)^2\pmod2$. By Lemma~\ref{lem:telescope},
$\bar S_i(\bar\omega)=\overline{(1-\omega)}^{-(i-1)}\bar A_i$; since
$\overline{1-\omega}=1+\bar\omega=\bar\omega^2$ (using $-1=1$, then $\bar\omega^2=\bar\omega+1$)
and $\bar\omega^3=1$, this is $\bar\omega^{-2(i-1)}\bar A_i=\bar\omega^{\,i-1}\bar A_i$. Squaring
gives the claim.
\end{proof}

\begin{lemma}[Congruence at the ramified prime]\label{lem:mod3}
$S_i(-1)\equiv(-1)^{i-1}A_i\pmod{\lambda}$.
\end{lemma}

\begin{proof}
Reduction $\Z[\omega]\to\Z[\omega]/\lambda\cong\F_3$ sends $\omega\mapsto-1$ (as
$\lambda=1+\omega$). Since $S_i\in\Z[q]$, evaluation commutes with this reduction:
$S_i(\omega)\bmod\lambda=S_i(-1)\bmod\lambda$. By Lemma~\ref{lem:telescope} the left side is
$(1-\omega)^{-(i-1)}A_i\equiv2^{-(i-1)}A_i\equiv(-1)^{i-1}A_i\pmod\lambda$ (using $2\equiv-1
\pmod3$, hence mod $\lambda$).
\end{proof}

\section{Proof of the main theorem}

\begin{theorem}[Shadow criteria]\label{thm:shadow}
For every $i\ge1$ and every integer sequence,
\[
\Phi_3(q)\mid S_i(q) \iff A_i\equiv0\pmod2, \qquad
\Phi_2(q)\mid S_i(q) \iff A_i\equiv0\pmod\lambda.
\]
\end{theorem}

\begin{proof}
By Lemma~\ref{lem:orbit}, $S_i(\zeta_3)$ is either $0$ or a unit of $\Z[\omega]$ (every nonzero
sixth root of unity is a unit); a unit never lies in the proper ideal $(2)$, so
$S_i(\zeta_3)=0\iff S_i(\zeta_3)\equiv0\pmod2$. By Lemma~\ref{lem:mod2} the right side holds iff
$\bar A_i^{\,2}=0$ in the field $\F_4$, i.e.\ iff $A_i\equiv0\pmod2$. This proves the first
equivalence. Likewise $S_i(-1)\in\{0,\pm1\}$ by Lemma~\ref{lem:orbit}, and $\pm1\not\equiv0
\pmod\lambda$ (as $\Z[\omega]/\lambda\cong\F_3$ and $\pm1\neq0$ there), so $S_i(-1)=0\iff
S_i(-1)\equiv0\pmod\lambda$, which by Lemma~\ref{lem:mod3} holds iff $A_i\equiv0\pmod\lambda$.
\end{proof}

\begin{proof}[Proof of Theorem~\ref{thm:main}]
Suppose $\Phi_6(q)\mid S_i(q)$. By Corollary~\ref{cor:level6}, $A_i=0$. Then trivially
$A_i\equiv0\pmod2$ and $A_i\equiv0\pmod\lambda$, so $\Phi_3(q)\mid S_i(q)$ and $\Phi_2(q)\mid
S_i(q)$ by Theorem~\ref{thm:shadow}. As $\Phi_2,\Phi_3,\Phi_6$ are pairwise coprime in $\Q[q]$,
their product $[6]_q$ divides $S_i(q)$. Evaluating at $q=1$ gives
$6=\Phi_2(1)\Phi_3(1)\Phi_6(1)\mid S_i(1)$.
\end{proof}

\begin{remark}
The three criteria together say more than Theorem~\ref{thm:main}: the entire divisibility
profile of $S_i$ at the divisors of $6$ is read off from the single Eisenstein integer $A_i$,
\[
\Phi_6\mid S_i\iff A_i=0,\qquad \Phi_3\mid S_i\iff2\mid A_i,\qquad \Phi_2\mid S_i\iff\lambda\mid A_i,
\]
so that, e.g., $\Phi_2\Phi_3\mid S_i$ without $\Phi_6\mid S_i$ exactly when $2\lambda\mid A_i$
but $A_i\neq0$. Note that $A_i$ is now built from \emph{suffix} partial sums of the digit
sequence $c_1,\dots,c_i$ (equivalently, prefix sums read from the last digit backward) rather
than prefix sums --- the direction forced on us by the fact that Morier-Genoud--Ovsienko's own
recursion \eqref{eq:mgo} looks one digit ahead for its quantum-integer coefficient. This is the
only respect in which the correct recursion changes the shape of the argument; the telescoping
mechanism itself, and both congruence lemmas, are otherwise identical to the single-digit case.
\end{remark}

\begin{thebibliography}{9}
\bibitem{MO20} S.~Morier-Genoud, V.~Ovsienko, \emph{$q$-deformed rationals and
$q$-continued fractions}, Forum Math.\ Sigma \textbf{8} (2020), e13.
\bibitem{KMRWY} T.~Kogiso, K.~Miyamoto, X.~Ren, M.~Wakui, K.~Yanagawa, \emph{Arithmetic on
$q$-deformed rational numbers}, Arnold Math.\ J.\ \textbf{11} (2025), no.~3.
\end{thebibliography}

\end{document}
